% Page budget: approximately 3--4 pages.
\section{Introduction}\label{sec:introduction}

Financial strategy research is cumulative. New signals, transformations,
execution rules, and risk controls produce candidate strategies; candidates
that survive verification enter a library and may contribute components to
later research. A successful research program therefore creates a growing
stock of executable policies, documented procedures, and reusable modules.
Growth is productive because it expands both the menu that can be operated
today and the material from which future strategies can be constructed.

Retention is not free. Every strategy must be stored, indexed, validated,
maintained, monitored, retrieved, documented, and governed. These activities
consume computing capacity and scarce human attention even when the strategy
is not currently selected. As libraries expand, the retention problem becomes
an economic decision rather than a housekeeping exercise: which strategies
can be removed, and which apparent redundancies are worth paying to keep?

This paper is organized around one question:
\begin{quote}
\emph{How much of a growing strategy library can be removed while preserving
both what the library can operate now and what it can help generate later?}
\end{quote}
The question separates two roles that are easy to conflate. A strategy is an
operating asset when its payoff profile improves the action available at a
current belief. It is a generative asset when it carries a signal,
transformation, execution component, risk control, or other module needed to
produce and verify a descendant. A strategy can matter through either role,
through both, or through neither.

Conventional current-performance pruning is tempting because the operating
role is directly visible. One can retain enough strategies to reproduce the
best current payoff at every belief and delete strategies that never attain
that envelope. Such a rule can be entirely correct as a statement about the
current action menu. It is incomplete as a statement about a research
library. A strategy that is dominated at every current belief may still be the
only carrier of a module used by a future project. Removing it changes what
can be generated even though it changes no payoff available today. Current
redundancy and generative redundancy are therefore different properties.

The distinction leads to a compact description of the library. For a raw
library \(L\), let \(F_L\) be its pointwise operating frontier: at each belief,
it records the highest current expected payoff among retained strategies. Let
\(C_L\) be its generative closure: the closed set of modules that can be
assembled from retained carriers. Alongside the filtered belief, the pair
\begin{equation}\label{eq:introduction-productive-state}
  K_L=(F_L,C_L)
\end{equation}
is the productive library state. The frontier records what can be operated
now; the closure records what the modeled research process can use to produce
later candidates. In the finite model, the raw generation, verification,
admission, and library-update process projects onto this pair. Projects may
have positive duration, and their terminal beliefs and admitted candidates may
be correlated; neither feature requires the raw strategy identifiers once the
frontier and closure are fixed.

This state reduction turns the organizing question into an exact
source-relative compression problem. Give each active strategy a positive,
fixed retention weight and write \(W(L)\) for total burden. For a specified
source library \(L\), exact safe compression solves
\begin{equation}\label{eq:introduction-safe-compression}
  \min_{L'\subseteq L} W(L')
  \quad\text{subject to}\quad
  F_{L'}=F_L,\qquad C_{L'}=C_L,
  \qquad s_0\in L'.
\end{equation}
The constraint is semantic rather than statistical: the retained sublibrary
must reproduce both productive objects exactly. Because the feasible family
is finite and contains the source, a minimum is attained. Every feasible
sublibrary has the same modeled current operating rewards, research menus,
candidate laws, future compressed updates, productive values, and available
optimal actions as the source. Compression changes the raw representation and
its burden, not the productive control process.

The exact constraint also gives a practical deletion certificate. Starting
from the source, delete a proposed strategy only when recomputing the library
shows that both its frontier and closure are unchanged. Rechecking after every
accepted deletion produces a certified safe path: each intermediate library
remains in the source's productive state. This certificate answers a local
feasibility question, however, not the global minimization in
\eqref{eq:introduction-safe-compression}. Earlier deletions can change which
remaining strategies are redundant. An endpoint at which no single further
deletion is safe can therefore have greater burden than another safe
sublibrary reached by a different deletion order or found by complete search.
Certified stepwise pruning and global minimum-resource compression are
distinct decisions.

A simple bridge illustrates why the closure constraint is economically
substantive. Suppose a strategy is dominated at every current belief but is
the unique carrier of a module required by a research project. If retained,
that project can combine the module with other capabilities, survive and pass
verification, and admit a valuable descendant after \(d\) periods. If the
carrier is deleted, the present operating frontier is unchanged, but the
project becomes unavailable and the descendant cannot be reached. When the
discounted, survival-adjusted descendant payoff exceeds research cost, that
difference is a positive delayed bridge margin. In the sharp construction,
frontier-only pruning removes the complete attainable margin, so normalized
loss equals one. The example does not say that every dominated strategy is
valuable; it shows why current-payoff preservation alone cannot certify that
a deletion is harmless.

Exact preservation is the primary problem. Two secondary variants ask what to
retain when scarcity is allowed to change productive opportunities. Over a
finite eligible catalog \(\mathfrak L_\theta\), a hard capacity \(B\) selects
the most productive library whose burden fits the constraint, while a resource
price \(\lambda\) selects value net of burden:
\begin{equation}\label{eq:introduction-resource-extensions}
\begin{aligned}
 V_\theta^\star(b;B)
   &=\max_{L'\in\mathfrak L_\theta:\,W(L')\le B}V_\theta(b,L'),\\
 P_\theta(b;\lambda)
   &=\max_{L'\in\mathfrak L_\theta}
      \{V_\theta(b,L')-\lambda W(L')\}.
\end{aligned}
\end{equation}
These are outer, one-time retention choices, not repeated eviction decisions
inside the Bellman problem. Their geometry is finite. Capacity value is
attained, monotone, and constant between attainable burdens, but complementary
modules can generate increasing unit increments. Penalized value is the
continuous convex piecewise-affine upper envelope of finitely many
value--burden lines, and selected burden falls weakly as the resource price
rises. Optimal library composition can change discretely; no concavity,
divisibility, or smooth policy path is imposed.

Productive value provides the link between exact compression and these
resource variants. The value of inserting a strategy separates into the
change in passive operating value and the change in the research-option
premium. Along a named parameter path and within a fixed selected policy, the
same accounting separates operational and generative sensitivity. Innovation
duration summarizes the discount exposure of future invention value. These
objects explain why a library is valuable without treating resource burden as
itself an operating or generative payoff.

The results are organized in four groups.
\begin{enumerate}
\item \emph{Raw-to-compressed state and exact safe compression.}
The finite raw research process projects onto belief, operating frontier, and
generative closure without changing productive values or available optimal
actions. Within a source library, a minimum-resource sublibrary preserving the
frontier and closure is attained
(\cref{thm:raw-to-compressed-projection,thm:minimum-safe-compression}).

\item \emph{Safe deletion, local versus global compression, and sharp loss.}
Frontier--closure equality characterizes safe deletion under the stated
detectability condition, and rechecking supplies a certified sequential path.
A deletion-irreducible endpoint need not minimize global burden, while
frontier-only pruning can attain the complete normalized delayed-bridge loss
(\cref{thm:local-global-safe-compression,%
thm:normalized-frontier-pruning-loss}).

\item \emph{Finite resource optimization and productive-value interpretation.}
Hard-capacity and resource-price selection have finite, potentially
complementary and discrete geometry. Operational and generative value explain
why the two safe-compression constraints matter, while named-path sensitivity
and fixed-exposure timing remain auxiliary comparative statics
(\cref{thm:capacity-value,thm:penalized-envelope,%
thm:value-decomposition,prop:channel-elasticity}).

\item \emph{Layered finite and financial evidence.}
Exact fixtures and complete small-library enumeration check identities and
boundary cases; registered randomized stress tests exercise the mechanisms in
frozen finite generators; and retrospective financial audits compare
stepwise-safe, HiGHS-reported minimum-resource, and frontier-only libraries
without
using held-out descendant quality in a pruning rule or objective
(\cref{sec:experiments}).
\end{enumerate}

The evidence levels are not interchangeable. Exact fixtures establish
machine-rechecked facts about specified finite instances, not universal
claims. Registered randomized results report frequencies for their frozen
synthetic generators, not population estimates. Mixed-integer solver status
is numerical global-search evidence and is reported separately from exact
rechecks of selected libraries, frontiers, closures, and burdens. The
financial exercises are retrospective information audits: structural and
validation inputs are fixed before held-out descendant quality is read, and
that quality is used only as an ex post opportunity diagnostic.

The scope is correspondingly narrow and explicit. States, beliefs, strategy
catalogs, project menus, and retained libraries are finite. Retention weights
are fixed and additive, and retention is chosen once outside the productive
dynamic program. The generator and verifier factor through the modeled belief
and frontier--closure state; provenance, age, multiplicity, time-varying
maintenance burdens, endogenous eviction, and endogenous verification are not
part of the primary model. Exact compression is source-relative and allows no
frontier or closure tolerance. Sensitivities are reported only on named paths
and, for elasticities, with positive denominators and a fixed selected policy;
finite changes replace derivatives at discrete changes.

The remainder of the paper follows this dependency order.
\Cref{sec:literature-boundary} fixes the boundary with partial-information
control, policy compression, resource-constrained selection, maintained
libraries, and generated financial strategies. \Cref{sec:model} defines the
raw process, the frontier--closure state, and retention burden.
\Cref{sec:compression} develops exact global and certified local compression;
\cref{sec:value} interprets productive value and finite capacity shadows;
\cref{sec:control} solves the canonical finite control and resource problems;
and \cref{sec:experiments} reports the controlled evidence. The primary
appendices contain proofs, core sensitivities, computational verification,
experiment-design information, and the verification-status summary. The
Online Supplement contains extended comparative statics, complete numerical
records, full randomized-study diagnostics, detailed financial-audit results,
and reproducibility records.

\clearpage
